\section{Independent reconstruction of the quaternary permanent certificates}
\label{sec:quaternary-reconstruction}

This section verifies inherited statements, rather than claiming priority
for their conclusions. The source locators are CUE\_GCT R-13,
R-24--R-27, in the corresponding numbered subsections of
\texttt{deliverables/CUE\_GCT.tex}. The checks below reconstruct the
matrices from their mathematical definitions. No archived matrix or
original checker is read. In fact, the supplied archive contains the
JSON specifications but not the referenced binary coefficient matrices
or C++ builders. This omission does not obstruct reconstruction.

\subsection{QC-01: the finite differential certificates}

\paragraph{Verification statement QC-01.}
For the polynomial map
\[
 \Phi_m:(\operatorname{Mat}_m\mathbb C)^4\longrightarrow
       \operatorname{Sym}^m(\mathbb C^4)^*,\qquad
 (A_0,A_1,A_2,A_3)\longmapsto
       \operatorname{per}(x_0A_0+x_1A_1+x_2A_2+x_3A_3),
\]
the source's certificates at $m=11,\ldots,17$ reproduce exactly.
Consequently
\[
 c_m=\binom{m+3}{3}\quad(11\leq m\leq16),\qquad c_{17}=1124.
\]
Here $c_m$ denotes the dimension of the image closure.

\begin{proof}
For each $m$, start at $s_0=20261121+1009m$ and set
\[
 s_{a+1}=1664525s_a+1013904223\pmod {2^{32}},\qquad
 c_a=(s_{a+1}\mathbin{\mathrm{div}}2^{16})\bmod101.
\]
These are integer coefficients in $\{0,\ldots,100\}$.
Set $L_{ij}(x)=\sum_{k=0}^3c_{4(mi+j)+k}x_k$ with zero-based
indices. Order degree-$m$ nonnegative quadruples in decreasing
lexicographic order. The evaluation of the differential has entries
\[
 E_{\alpha,4(mi+j)+k}=
 \alpha_k\operatorname{per}(L(\alpha)_{\widehat i,\widehat j}).
 \tag{QC-01.1}
\]
Evaluation is a linear map on the coefficient target and cannot
increase rank.

For completeness, the full arithmetic certificate has a short
definition. If $A$ is a numerical $m\times m$ matrix, set
\[
 F(\varnothing)=1,\qquad
 F(S)=\sum_{j\in S}A_{|S|-1,j}F(S\setminus\{j\}).
 \tag{QC-01.2}
\]
The induction on $|S|$ identifies $F(S)$ with the permanent on
the first $|S|$ rows and columns $S$. Reverse differentiation starts
with $B([m])=1$ and otherwise zero. For decreasing subsets $S$ and
each $j\in S$, it adds
\[
 B(S)A_{|S|-1,j}\ \text{to }B(S\setminus\{j\}),\qquad
 B(S)F(S\setminus\{j\})\ \text{to }G_{|S|-1,j}.
\]
The chain rule gives $G_{ij}=\partial\operatorname{per}(A)/\partial
A_{ij}$, exactly the complementary permanent in (QC-01.1).
Thus these recurrences define every entry without trusting a stored
array.

The standard-library checker \texttt{check\_projection.py}
reconstructed the full matrices and obtained
\[
\begin{array}{c|r|r|r}
 m&\text{matrix rows}&\text{certified rank}&
           \text{selected determinant modulo }101\\\hline
11&364&364&42\\
12&455&455&87\\
13&560&560&25\\
14&680&680&69\\
15&816&816&73\\
16&969&969&8\\
17&1140&1124&88
\end{array}
\]
For $m\leq16$, columns are selected greedily from left to right,
with the first available nonzero pivot row. The determinant uses
the original row order. For $m=17$, use the first $1124$ rows and
omit exactly the columns
\[
 \{68i+67:1\leq i\leq15\}\ \cup\
 \{1091+4j:0\leq j\leq16\}.
\]
The latter matrix is square. The checker gives complete selected
column lists and SHA-256 digests of all rebuilt entries; these
digests document the run and are not substitutes for reconstruction.

The Gaussian elimination is exact. Its packed implementation is
just scalar modular elimination with independent integer lanes.
For $b=\lceil\log_2(p+1)\rceil$ and $0\leq z<p^2$, put
\[
 q=\left\lfloor {z\lfloor2^{2b}/p\rfloor\over2^{2b}}\right\rfloor.
\]
Then $q$ is $\lfloor z/p\rfloor$ or one less, because the difference
before flooring is less than $p^2/2^{2b}<1$. One correction therefore
reduces $z-pq$ modulo $p$. Intermediate products use fewer than $3b$
bits; the implementation reserves at least $4b$ bits per lane.
The fixed positive offset in subtraction prevents a borrow between
lanes. Hence it computes precisely ordinary field elimination.
An additional independent scalar implementation agreed on rectangular
test matrices, and direct permutation sums checked the gradient
recurrence at sizes through six. The main certificate is nevertheless
the exact reconstructed matrix and its nonzero minor, not those tests.

A nonzero residue is a nonzero integer minor and hence a nonzero
minor over $\mathbb C$. The ranks in the table imply the lower
bounds stated. For the upper bound at $17$, the row--column torus
\[
 H_m=\{(u,v)\in(\mathbb C^*)^{2m}:\prod_i u_i\prod_jv_j=1\}
\]
preserves the permanent of the pencil under $L_{ij}\mapsto u_iv_jL_{ij}$.
Its generic stabilizer consists of $(u_i,v_j)=(t,t^{-1})$, because
all entries are nonzero linear forms. The generic fiber has dimension
at least $2m-2$, giving $c_m\leq4m^2-2m+2$.
At $m=17$ this upper bound is $1124$; the other upper bounds are
the target dimensions. This proves QC-01.
\end{proof}

\subsection{QC-02: the bordered induction and $c_{18}$}

\paragraph{Verification statement QC-02.}
The conditional induction of CUE\_GCT R-25 is valid. Its finite
$m=17$ certificate also reproduces, so $c_{18}=1262$.

\begin{proof}
Let $A$ be an $m\times m$ matrix of ternary linear forms and let
$u,v$ be vectors of ternary linear forms. Write
\[
 P=\operatorname{per}(A),\quad
 Q=DP_A[uv^{\mathsf T}],\quad
 R=\tfrac12D^2P_A[uv^{\mathsf T},uv^{\mathsf T}],
\]
and let $a_i,b_j,t_i,t'_j$ be respectively the derivatives of $Q,R$
in the auxiliary scalar coordinates $u_i,v_j$, followed by substitution
of the indicated linear forms. With $E^*=\langle x,y,z\rangle$, set
\[
 C_A(u,v)=E^*P+\sum_iE^*a_i+\sum_jE^*b_j+
                   \langle t_i,t'_j\rangle
       \subseteq\operatorname{Sym}^{m+1}E^*.
 \tag{QC-02.1}
\]
There are $8m+3$ displayed columns. The identity
$\sum_i u_i a_i=\sum_jv_jb_j$ is a nonzero constant-coefficient
relation among the linear-multiple columns whenever $u$ or $v$
is nonzero. Thus their span has dimension at most $8m+2$.

Assume that equality holds for some ternary data. Add a fourth
variable $\ell$. On the irreducible affine space of quaternary data,
the ternary rank condition after $\ell=0$ is a nonempty open
condition. The condition that the quaternary differential at $A$
have rank $c_m$ is another nonempty open condition. They hold
simultaneously somewhere. At such a point form
\[
 B_\varepsilon=\begin{pmatrix}A&\varepsilon u\\
                  \varepsilon v^{\mathsf T}&\ell\end{pmatrix}.
\]
Every permutation uses either the bottom-right entry or both
off-diagonal blocks, so exactly
\[
 \operatorname{per}(B_\varepsilon)=\ell P+\varepsilon^2Q.
 \tag{QC-02.2}
\]
Bottom-right variations produce $\langle x,y,z,\ell\rangle P$.
Off-diagonal variations, divided by $\varepsilon$, produce all
linear multiples of $a_i,b_j$ for $\varepsilon\ne0$.
For a constant vector $h$, vary the upper-right block by $\ell h$
and the upper-left block by $-\varepsilon hv^{\mathsf T}$.
The order-one terms cancel, and (QC-02.2) gives the exact derivative
\[
 -\varepsilon^3D^2P_A[hv^{\mathsf T},uv^{\mathsf T}]
             =-\varepsilon^3\sum_i h_i t_i.
\]
Varying the lower-left block similarly produces the $t'_j$.
All these variations still have linear-form entries and belong to
the differential of the same pencil, not different deformations.

Choose $8m+2$ such fixed columns whose restrictions to $\ell=0$
are independent. Choose $c_m$ upper-left variations $U_i$ with
$DP_A[U_i]$ independent. Their derivatives at $B_\varepsilon$ are
$\ell DP_A[U_i]+\varepsilon^2D_AQ[U_i]$.
At $\varepsilon=0$ the combined columns are independent: restriction
to $\ell=0$ first removes the normal columns' coefficients, and
multiplication by $\ell$ is injective on the remaining span.
Their nonzero limiting minor is therefore nonzero for some
$\varepsilon\ne0$. It follows that
\[
 c_{m+1}\geq c_m+8m+2.\tag{QC-02.3}
\]

For the finite antecedent at $m=17$, use $p=1021$, $s_0=20330804$,
the same linear congruential recurrence as QC-01, and coefficients
$c_a=(s_{a+1}\mathbin{\mathrm{div}}2^8)\bmod1021$.
The first $3\cdot17^2$ coefficients define $A$ in row, column,
variable order, followed by $3\cdot17$ coefficients each for $u,v$.
Evaluate (QC-02.1) at the $190$ degree-$18$ nonnegative triples,
in decreasing lexicographic order. The column order is
\[
 x_kP;\quad x_ka_i;\quad x_kb_j;\quad t_i;\quad t'_j,
\]
with index first and variable second. On the first $138$ rows,
omit column $104$. The resulting determinant is $501\pmod{1021}$.

The new checker \texttt{check\_border18.py} obtains all $26{,}410$
entries directly. Apply (QC-01.2) to $A+tuv^{\mathsf T}$ modulo
$t^2$ and reverse differentiate that truncated recurrence. If
$G(t)=G_0+tG_1+O(t^2)$ is its matrix-entry gradient, the chain rule
gives
\[
 a=G_0v,\qquad b=G_0^{\mathsf T}u,\qquad
 t=G_1v,\qquad t'=G_1^{\mathsf T}u.
\]
For example $\partial_{u_i}P(A+tuv^{\mathsf T})=
t\sum_jv_jG_{ij}(t)$, whose $t^2$ coefficient is $\partial_{u_i}R$.
This proves that the checker reconstructs exactly the required
columns. The nonzero minor gives dimension at least $138$, and
the universal relation gives at most $138$. Applying (QC-02.3)
and QC-01 yields $c_{18}\geq1124+138=1262$. The torus upper
bound is $1262$, proving equality.
\end{proof}

\subsection{QC-03: finite cores imply the uniform endpoint rank}

\paragraph{Verification statement QC-03.}
For every $m\geq17$ there are ternary data with
$\dim C_A(u,v)=8m+2$. Consequently
\[
 c_m=4m^2-2m+2\quad(m\geq17).
 \tag{QC-03.1}
\]
Together with the settled sizes $m\leq10$ and QC-01, this verifies
the dimension assertion of CUE\_GCT R-27 for every size. We give
the insertion proof, because a list of finite ranks by itself
does not establish a uniform assertion.

\begin{proof}
Put $L=m-1$, $h=yz$, and take the Jacobi pencil
\[
 A_{ii}=\alpha_i x,\quad A_{i,i+1}=y,\quad A_{i+1,i}=\delta_i z.
\]
For an interval let $P_{a,b}$ be its permanent, with
$P_{a,a-1}=1$. A path matching either fixes its last vertex or
pairs it with its predecessor, giving
\[
 P_{a,b}=\alpha_bxP_{a,b-1}+\delta_{b-1}hP_{a,b-2}.
 \tag{QC-03.2}
\]
Set $D_i=\prod_{e=i}^{L-1}\delta_e$ and
$E_i=\prod_{e=0}^{i-1}\delta_e$. Let $u,v$ have endpoint entries
$(U_0,U_1)$ and $(V_0,V_1)$, respectively. Forced paths in
complementary matchings give
\begin{align*}
 a_i&=V_0y^iP_{i+1,L}+V_1D_i z^{L-i}P_{0,i-1},\\
 b_i&=U_0E_i z^iP_{i+1,L}+U_1y^{L-i}P_{0,i-1},\\
 t_i&=2V_0V_1\bigl(U_0D_i z^{L-i}P_{1,i-1}
                       +U_1y^iP_{i+1,L-1}\bigr),\\
 t'_i&=2U_0U_1\bigl(V_0y^{L-i}P_{1,i-1}
                       +V_1E_i z^iP_{i+1,L-1}\bigr).
\end{align*}
In the last two lines omit the first summand at $i=0$ and the
second at $i=L$. Each path fixes all entries between the removed
row and column; outside it the remaining matchings are interval
matchings. For the second derivatives there are two assignments
of the rank-one entries, which accounts for the factor two.
This proves the formulas rather than assuming them from a numerical
check.

Let $C(\varepsilon)$ be the coefficient matrix of (QC-02.1) along
\[
 (U_0,U_1,V_0,V_1)=(x,y,x,z)+
       \varepsilon(y+2z,3x+4z,5y+6z,7x+8y).
 \tag{QC-03.3}
\]
Its entries are polynomials of degree at most three in $\varepsilon$.
At zero the cyclic action $(x,y,z)\mapsto(x,\omega y,\omega^{-1}z)$,
$\omega^m=1$, decomposes the columns by character. The columns
$a_i,t_i$ have character $i$, and $b_i,t'_i$ have character $-i$.
Characters $0,1,-1$ have nine columns each; every other character
has eight.

\paragraph{Ordinary characters: the finite core calculation.}
For $2\leq r\leq m-2$, put $j=m-r$ and set $\alpha_i=1$ for
this calculation. In character $r$ use the columns
\[
 xa_r,ya_{r-1},za_{r+1},xb_j,yb_{j+1},zb_{j-1},t_r,t'_j.
\]
They have the form $y^rp+z^jD_rq$. On the nonzero-edge locus,
divide the second component by $D_r$ and set $K=E_j/D_r$.
The resulting pairs are
\[
\begin{array}{c|c|c}
 &p&q\\\hline
1&x^2P_{r+1,L}&xP_{0,r-1}\\
2&xP_{r,L}&\delta_{r-1}hP_{0,r-2}\\
3&xhP_{r+2,L}&\delta_r^{-1}P_{0,r}\\
4&xP_{0,j-1}&Kx^2P_{j+1,L}\\
5&P_{0,j}&K\delta_jxhP_{j+2,L}\\
6&hP_{0,j-2}&K\delta_{j-1}^{-1}xP_{j,L}\\
7&2xhP_{r+1,L-1}&2x^2P_{1,r-1}\\
8&2x^2P_{1,j-1}&2KxhP_{j+1,L-1}
\end{array}\tag{QC-03.4}
\]
The scalar $K$ reduces to
\[
 K={\prod_{e=0}^{\min(r,j)-1}\delta_e\over
             \prod_{e=\max(r,j)}^{m-2}\delta_e}.
 \tag{QC-03.5}
\]
To specify the cores, initially set
$\delta_i=7i^2+13i+36\pmod{101}$, replacing a zero by one.
For $(r,j)=(6+\max(-g,0),6+\max(g,0))$, $-4\leq g\leq4$,
set $\delta_2=\delta_{m-4}=0$ and take $K=14$; also set
$\delta_{\min(r,j)+1}=0$ when $|g|=4$.
For $r+j=17$ with $\min(r,j)=2,3,4,5$, set only
$\delta_{\min(r,j)+1}=0$ and define $K$ by (QC-03.5).
There are seventeen cores.

The rows of (QC-03.4) are the coefficients of
$x^{j+1-2k}h^k$, followed by those of $x^{r+1-2k}h^k$, in
increasing $k$. The checker \texttt{check\_ordinary.py} explicitly
enumerates nonadjacent edge sets of each interval and obtains
the following $8\times8$ determinants in the ordered row lists:
\[
\begin{array}{r r|l|r}
r&j&\text{ordered rows}&\det\pmod{101}\\\hline
10&6&0,1,4,2,3,5,6,7&73\\
9&6&0,1,4,2,3,5,6,7&58\\
8&6&0,1,4,2,3,5,6,7&16\\
7&6&0,1,4,2,3,5,6,7&8\\
6&6&0,1,4,2,3,5,6,7&53\\
6&7&0,1,5,3,2,6,7,8&51\\
6&8&0,1,5,3,4,6,2,7&72\\
6&9&0,1,6,3,4,7,2,8&20\\
6&10&0,1,6,3,4,7,2,8&90\\
2&15&0,1,9,3,4,10,6,7&51\\
15&2&0,1,2,3,4,5,6,7&16\\
3&14&0,1,8,3,4,9,6,7&88\\
14&3&0,1,3,4,2,5,6,7&93\\
4&13&0,1,8,3,4,9,6,7&23\\
13&4&0,1,3,2,4,5,6,7&24\\
5&12&0,1,7,3,4,8,6,2&83\\
12&5&0,1,4,2,3,5,6,7&40
\end{array}
\]
The last eight $K$ residues are, in pairs, $34,61,88,82$.
All denominators in the selected matrices are nonzero modulo $101$.

These normalized cut data are limits of nonzero-edge data.
A zero edge in the middle interval absent from (QC-03.5) does not
affect $K$. If outer zero edges occur in equal numbers on the two
sides, replace them by $\tau$ times nonzero constants. Their
powers of $\tau$ cancel in (QC-03.5); changing one left constant
realizes any prescribed nonzero $K$. Hence every nonzero cut
minor proves a nonzero rational minor on the nonzero-edge locus.

There are two exact insertions. When $r,j\geq6$, insert a vertex
with diagonal $x$ at each of the cuts $2,m-4$, replacing each cut
edge by two zero edges. Every interval in either component of
(QC-03.4) gains one isolated vertex, so $(p,q)$ becomes $(xp,xq)$
and $(m,r,j)$ becomes $(m+2,r+1,j+1)$. The outer zero counts
remain equal, so $K$ can remain fixed. When $j-r\geq4$, insert
one vertex at the cut $r+1$. Every $p$ interval gains that vertex,
and every $q$ interval avoids it; the map is $(p,q)\mapsto(xp,q)$,
and $(m,r,j)\mapsto(m+1,r,j+1)$. The case $r-j\geq4$ is
$(p,q)\mapsto(p,xq)$. These assertions follow by checking the
endpoints of the eight intervals in each column of (QC-03.4),
including the shifted intervals with $r+2$ or $j-2$.
Multiplication by $x$ is injective, so the ranks are preserved.
One step of every applicable insertion was also checked entry by
entry by the independent program.

For any $r,j\geq6$, choose the core with the same signed gap,
truncated to $[-4,4]$. Paired insertions raise the smaller index
to its target; if necessary central insertions then raise the
larger index. If the smaller index is $2,3,4,5$ and $r+j\geq17$,
the corresponding size-$17$ core and central insertions suffice.
Thus every ordinary block has rank eight on a nonempty open
subset of the actual Jacobi parameter space.

\paragraph{The boundary certificate and its insertion.}
Retain every row and every base-character column with character
in $B_m=\{0,\pm1,\ldots,\pm5\}\pmod m$. Call the projected
matrix $S(\varepsilon)$; it has $91$ columns. At size $17$ take
\begin{align*}
\alpha={}&(86,18,1,17,95,90,80,90,48,100,9,60,2,25,78,53,54),\\
\delta={}&(34,1,37,24,47,1,59,41,0,92,38,56,75,6,44,14).
\end{align*}
Order the degree-$18$ monomials by decreasing $x$ exponent and
then decreasing $y$ exponent. Order columns as $xP,yP,zP$,
then $xa_i,ya_i,za_i$, then $xb_i,yb_i,zb_i$, then $t_i,t'_i$.
Retain the indicated complete character blocks.

The zero edge at $8$ is interpreted after row normalization:
for negative extreme rows divide by $D=\prod\delta_i$ before
taking the cut limit. Every term in such a retained row has the
long lower path, and its residual interval avoids the cut. Thus
the cut factor cancels. Computationally, set the cut to one in
these rows and divide by the product of the other edges; in the
other rows set the cut to zero. The other-edge product is one
modulo $101$ for the displayed data. This is a rational limit of
invertible row scalings, not the literal degenerate-pencil rank.

Write $S=\sum_{i=0}^3\varepsilon^iS_i$ and form the lower
triangular block matrices $T_s=(S_{i-j})_{0\leq i,j\leq s}$,
with negative subscripts zero. Direct sparse-polynomial expansion
in \texttt{check\_endpoint.py} yields
\[
\begin{array}{c|c|r|r}
 &\text{size}&\text{rank}&\text{ordered-pivot minor}\pmod{101}\\\hline
T_0&124\times91&84&99\\
T_1&248\times182&172&74\\
T_2&372\times273&262&45
\end{array}\tag{QC-03.6}
\]
The pivot rule is: scan columns left to right, choose the first
unused nonzero row, and eliminate below after normalization.
The selected original row labels are ordered by pivot time.
In this ordered convention the determinant is the product of
the pre-normalization pivots. This specifies the minors fully.

Insert an isolated vertex of diagonal $x$ at the central cut,
replacing the zero edge by two zero edges. Repeat at the leftmost
central cut. For retained first-derivative columns, the endpoint
distance is at most six; for second-derivative columns it is at
most five. In each formula above, a short path has a long residual
interval, while a long path has a short residual interval.
The original cut is between vertices eight and nine. Every long
residual interval contains each inserted vertex; every short
interval avoids it. Every long forced path crosses all central
cuts; every short path avoids them.

The retained integer-weight windows are precisely
\[
 [-5,5],\qquad[m-5,m+1],\qquad[-m-1,-m+5].
\]
A short-path term has absolute weight at most eight, and a long-path
term at least $m-9$. For $m\geq17$ the central retained window
therefore contains only short paths, and the extreme retained
windows only long paths. In the three windows the insertion is
respectively the injective monomial map
\[
 f\longmapsto xf,\qquad f\longmapsto yf,\qquad f\longmapsto zf.
 \tag{QC-03.7}
\]
The negative long path gains new edge factors, all canceled by
the normalization by $D$. Its residual interval avoids those
edges. Thus (QC-03.7) holds for the entire polynomial matrix,
including every coefficient of $\varepsilon$. New rows outside
these images are zero. It preserves every $T_s$ rank.
The checker separately verified the entire embedding at sizes
$18,19,20,21$, including order three; the preceding argument, not
those four tests, proves the assertion for all sizes.

Replacing all central zeros by $\tau$ times nonzero constants
realizes the normalized cut as a finite rational limit of actual
nonzero-edge pencils. The nonzero limit minors in (QC-03.6)
therefore define nonempty rank opens at every size.

\paragraph{The seven base relations and first normal bound.}
At the base endpoints $(x,y,x,z)$ there are seven independent
relations. Five of them are
\begin{align*}
&xa_0+ya_L-xb_0-zb_L=0,\\
K_{00}:={}&yb_1-\delta_0za_1+\alpha_0xa_0-\alpha_0xb_0=0,\\
K_{01}:={}&t_0-t'_0=0,\\
K_{10}:={}&\delta_0t_1-2xb_L+\alpha_0t'_L=0,\\
K_{11}:={}&\alpha_L\delta_0
       (\delta_1za_2-ya_0+\alpha_1xa_1)\\
&\hspace{7mm}-\alpha_0\alpha_1
       (\delta_{L-1}zb_{L-1}-yP+\alpha_Lxb_L)=0.
\end{align*}
The last two are the negative-character transforms of $K_{10},K_{11}$.
Explicitly, transpose $A$, interchange $y,z$, and conjugate by
$\operatorname{diag}(1/E_i)$. The pencil returns to Jacobi form
and its endpoints become $(x,Dy,x,D^{-1}z)$.
Next apply $(y,z)\mapsto(ty,t^{-1}z)$ and conjugate by
$\operatorname{diag}(t^{-i})$, where $t^m=D^{-1}$.
This restores the base endpoints and exchanges characters $r,-r$.
It is invertible over $\mathbb C$ and takes arbitrary endpoint
perturbations to arbitrary endpoint perturbations.

The relations follow directly from (QC-03.2). For example the two
parentheses in $K_{11}$ at the base are
$2\alpha_0\alpha_1D_1x^2z^{m-1}$ and
$2\alpha_LDx^2z^{m-1}$, respectively. Independence follows by
characters: in character zero the Hessian support distinguishes
$K_{01}$, and the index-one columns distinguish $K_{00}$ from
scaling; in each of characters $\pm1$ the Hessian support
distinguishes the first relation from the second, which has a
nonzero $yP$ or $zP$ coefficient. Thus $\operatorname{rank}C_0
\leq8m-4$ on the nonzero-parameter locus.

The boundary minor has rank at least $84$, and every character
outside $B_m$ has eight independent columns on a nonempty open
by the ordinary-core argument. Intersect these finitely many opens.
They lie in an irreducible parameter space, so the intersection
is nonempty and
\[
 \operatorname{rank}C_0=84+8(m-11)=8m-4.
 \tag{QC-03.8}
\]
Its kernel is exactly the seven relations just given, wholly in
characters $0,\pm1$; the projected kernel is the same kernel.
This also avoids any need to import the separate exceptional-
character calculation from R-26.

For arbitrary endpoint forms the same recurrences give
\[
 K_{11}=2\alpha_0\alpha_1\alpha_LD\,
    xz^{L-1}(xV_1-zU_0)-\alpha_0\alpha_1(U_1-y)P.
 \tag{QC-03.9}
\]
Its first derivative is in the base image. Here are explicit
witnesses. For $r=2,3$, set $j=m-r$ and put
\begin{align*}
 R_A(r)&=ya_{r-1}-\alpha_rxa_r-\delta_rza_{r+1}
        =-2\alpha_rD_r xz^{m-r}P_{0,r-1},\\
 R_B(r)&=yb_{j+1}-\alpha_jxb_j-\delta_{j-1}zb_{j-1}
        =-2\alpha_jE_j x^2z^jP_{j+1,L}.
\end{align*}
At $r=2$ they give $x^3z^{m-2}$ and $xyz^{m-1}$.
At $r=3$ they give $x^2yz^{m-2}$ provided
\[
 \Delta=\alpha_0\alpha_1\alpha_2\delta_{L-1}
       -\alpha_{L-1}\alpha_L(\alpha_2\delta_0+\alpha_0\delta_1)
       \ne0.
\]
This is a nonempty open condition (all parameters one give $-1$).
The identities
\begin{align*}
\delta_1za_2-ya_0+\alpha_1xa_1
 &=2\alpha_0\alpha_1D_1x^2z^{m-1},\\
\delta_0za_1-xP+\alpha_0xa_0&=2\alpha_0D xz^m
\end{align*}
give the other two possible monomials. Together with linear
multiples of $P$, these are all the terms in the first derivative
of (QC-03.9). The negative-character transform gives another
independent kernel vector with zero first normal. Scaling does
as well, because
$U_0a_0+U_1a_L-V_0b_0-V_1b_L=0$ holds along the entire family.
Therefore the first normal map
\[
 N_1:\ker C_0\longrightarrow\operatorname{coker}C_0,
             \quad k\longmapsto[C_1k]
\]
has rank at most four on a nonempty open set.

\paragraph{Second normals and the conclusion.}
For any polynomial matrix $M=M_0+\varepsilon M_1+\cdots$, let
$r_0=\operatorname{rank}M_0$ and $r_1=\operatorname{rank}N_1$.
For $k_0\in\ker N_1$, choose $k_1$ with
$M_0k_1=-M_1k_0$ and define
$N_2(k_0)=[M_2k_0+M_1k_1]$ modulo the base image and first-normal
image. This is independent of the choice, since the difference
of two choices belongs to $\ker M_0$. Write $r_2=\operatorname{rank}N_2$.
Successively eliminating a unit block of $M_0$ and an order-one
block gives
\[
 \operatorname{rank}T_0=r_0,\quad
 \operatorname{rank}T_1=2r_0+r_1,\quad
 \operatorname{rank}T_2=3r_0+2r_1+r_2,
\]
and $\operatorname{rank}_{\mathbb C(\varepsilon)}M\geq r_0+r_1+r_2$.
Equivalently, modulo $\varepsilon^{s+1}$ a pivot of orders
$0,1,2$ contributes $s+1,s,s-1$ to the Toeplitz rank.

These normal maps for the projected boundary matrix are the full
normal maps. Indeed, replacing one endpoint monomial shifts its
character by at most two. The base kernel has characters $0,\pm1$,
so $C_1k_0$ has characters at most $\pm3$. Any correction can
be projected onto those complete base blocks. The second residual
$C_2k_0+C_1k_1$ then has characters at most $\pm5$.
Outside blocks cannot alter either quotient. Retaining complete
blocks is essential for this assertion.

Intersect all ordinary-block, boundary-minor, and $\Delta$ opens
at one actual pencil. For $S$ we have $r_0=84$, $r_1\leq4$,
and (QC-03.6). Hence $172\leq168+r_1$ forces $r_1=4$, and
$262\leq252+8+r_2$ gives $r_2\geq2$. Together with (QC-03.8),
the full matrix has generic rank at least $(8m-4)+4+2=8m+2$.
The universal scaling relation gives the reverse bound.
A nonzero maximal minor is a nonzero polynomial in $\varepsilon$,
so some actual value attains that rank. This proves the antecedent
of (QC-02.3) for every $m\geq17$. Starting with $c_{17}=1124$,
the increment $8m+2$ is exactly the increment of $4m^2-2m+2$.
The torus upper bound proves (QC-03.1) by induction.
\end{proof}

\subsection{QC-04: the indicated derived-restriction degree}

\paragraph{Verification statement QC-04.}
Let $R=\mathbb C[X_{ij}]$, let $J_m$ be the ideal of submaximal
subpermanents, and let $S$ be a general four-variable linear
quotient of $R$. For $m\geq17$,
\[
 \operatorname{Tor}_1^R(R/J_m,S)_m=0.
\]
This is only the stated internal degree, not full Tor independence.

\begin{proof}
The full-variable input is Efremenko--Landsberg--Schenck--Weyman,
\emph{On minimal free resolutions of sub-permanents and other ideals
arising in complexity theory}, author version
\texttt{ELSWmfr2-29.pdf}, Theorem 1.1, printed page 2.
Putting $\kappa=m-1$ and their $j=2$ gives exactly $2m-2$
linear first syzygies. They are the differences of row Laplace
expansions and differences of column Laplace expansions.
Their independence after a general restriction is elementary:
a constant combination is the vector
$((\alpha_i+\beta_j)L_{ij})_{i,j}$ with
$\sum\alpha_i=\sum\beta_j=0$. When every $L_{ij}$ is nonzero,
vanishing forces all $\alpha_i$ to be equal and all $\beta_j$
their negatives; the two zero sums give zero coefficients.

Tensoring a minimal free resolution with $S$ and taking degree
$m$ gives the relevant complex
\[
 \mathbb C^{2m-2}\longrightarrow S_1^{m^2}
                    \longrightarrow S_m.
\]
The first map is injective generically. The second is the
permanent coefficient-map differential and has generic rank $c_m$
in characteristic zero. Thus
\[
 \dim\operatorname{Tor}_1^R(R/J_m,S)_m
       =4m^2-(2m-2)-c_m.
\]
QC-03 makes the right side zero. No assertion about other degrees
or about every positive characteristic is implied.
\end{proof}

\subsection{Geometric dictionary and limits of the conclusion}

The coefficient maps are morphisms of affine schemes defined by
integer polynomials. Each displayed minor is a section of an
exterior power of their tangent map. Its nonvanishing defines a
principal open; the arithmetic calculations prove that the
corresponding characteristic-zero open is nonempty. The torus
argument is an assertion about the fibers of a group-scheme action.
The endpoint argument works with maps of free modules over a
power-series ring: its successive normal maps measure ranks after
orders zero, one, and two. The insertion is a monomial injection
of the relevant graded modules, not an extrapolation of dimensions.

For QC-04 the coordinate complex of the derived intersection of
$\operatorname{Spec}(R/J_m)$ with the general four-plane is
$(R/J_m)\otimes_R^{\mathbf L}S$; QC-04 kills its first homology
in one internal degree. Framework locators, carried from the
supplied CUE\_GCT dictionary, are SGA 3, Expos\'e I,
Propositions 4.7.2--4.7.3 (actions and gradings), SGA 6,
Expos\'e III, Definition 1.5 and Notation 1.6 (Tor independence
and derived tensor product), and EGA II, Proposition 2.5.2
(graded modules and associated projective sheaves). The archive
transcription at cote 133, scan 028, lines 6--17 supplies the
tensor/symmetric-operation vocabulary; scan 019, lines 45--56
distinguishes an equivariant class from its augmentation. These
framework references are not claims that the cited texts contain
the modern permanent-dimension theorem. No motive, purity, or
$\ell$-adic assertion is needed for the present proofs.

At $(m,n)=(17,23)$ the verified affine dimensions are $1127$ for
the padded four-variable image and $1060$ for the determinant
image. At $(8,9)$ the inherited $c_8=165$ similarly gives $168$
and $164$. For the determinant dimension, a primary literature
locator is Reichstein--Vistoli, \emph{On the dimension of the locus
of determinantal hypersurfaces}, Canadian Mathematical Bulletin
60 (2017), 613--630, Theorem 1.3; Leal--Lozano Huerta--Vite,
\emph{The Noether--Lefschetz locus of surfaces in $\mathbb P^3$
formed by determinantal surfaces}, arXiv:2303.09028v3,
Corollary 3.1, gives the same projective dimension $2n^2+1$.
The latter's Corollary 1 identifies the smooth determinantal
loci with open parts of Noether--Lefschetz components.

None of these dimension verifications computes a highest-weight
equation or either complete isotypic multiplicity. In particular,
they do not supply the explicit obstruction requested in item 1.
The qualitative strict-direction argument of CUE\_GCT G-17/G-18
was read separately: its quotient dimensions explain eventual
growth, and its period argument removes a degree congruence.
For $n>\sqrt2m$ it gives only eventual determinant dominance
in each fixed four-row direction; it leaves finitely many early
degrees undecided. The enormous effective bounds in those
statements are not used or newly certified here.
